\lesson{4}{Sep 22 2021 Wen (12:43:32)}{Complex Numbers}{Unit 1}

\begin{definition}[Imaginary Numbers]
  \label{def:imaginary_numbers}

  Here's a quick question. Solve this equation: $a \times b = -1$. They must not
  be identical factors, meaning it cannot be $1 \times -1 = -1$.

  To solve this problem, we came up with the imaginary number $i$. It's defined
  to:

  \[ i = \sqrt{-1} \].

  With this definition, the square root of a negative radicand, in addition to
  positive radicands can be simplified.
\end{definition}

\subsubsection{History of the Imaginary Numbers}
\label{sub_sub_sec:history_of_the_imaginary_numbers}

Thought history, there have been times when a mathematical construct was
invented before its purpose was discovered. For example, early man knew the
concept of $1$ before knowing $0$. Numbers, at that time, were used to count
sheep and other animals as a source for food. There was no need to define the
concept of $0$. But, when fractions were defined, they were thought of useless.
What purpose could there be to have a number between $0$ and $1$. But today,
fractions become part of our daily life. Imaginary numbers have the same
history. It's use will come out eventually, we just don't know when.

Solving an equation such as: $x^{2} - 9 = 0$ should be already familiar. The
goal is to isolate the variable on one side by first adding $9$ to both sides.

\begin{equation}
  \label{eqt:history_of_the_imaginary_numbers_1}
  \begin{split}
    x^{2} &= 9 \\
    \sqrt{x^{2}} &= \sqrt{9} \\
    \boxed{x = \pm 3}
  \end{split}
.\end{equation}

Notice that the $\sqrt{9}$ is indicated as  $\pm 3$. The reason this is because
there are two pairs of identical factors that produce $9$ when multiplied: $3
\times 3 = 9$ and $-3 \times -3 = 9$.

But, how do you solve an equation such as: $x^{2} + 1 = 0$? You begin just as in
the last equation. Isolate the variable, but this time, subtract $1$ from both
sides of the equal sign.

\begin{equation}
  \label{eqt:history_of_the_imaginary_numbers_2}
  \begin{split}
    x^{2} + 1 &= 0 \\
    x^{2} &= -1 \\
    \sqrt{x^{2}} &= \sqrt{-1} \\
  \end{split}
.\end{equation}

This is what troubled mathematicians for centuries. All of the radicans they
have been working with have all been positive. So, to overcome this obstacle,
they invented $i$.

\begin{example}
  \label{exm:history_of_the_imaginary_numbers_3}

  Let's try and simplify the following expression: $\sqrt{-4}$
  
  \begin{equation}
    \label{eqt:history_of_the_imaginary_numbers_4}
    \begin{split}
      \sqrt{-4} &= \sqrt{-1} \times \sqrt{4} \\
            &= i \times \sqrt{4} \\
            &= i\sqrt{4} \\
            &= \boxed{2i}
    \end{split}
  .\end{equation}
\end{example}

% subsubsection history_of_the_imaginary_numbers (end)

\subsubsection{Using $i$ to help solve equations}
\label{sub_sub_sec:using_i_to_help_solve_equations}

You already know that $i = \sqrt{-1}$, but what about the different powers of
$i$. Take a look at the table below (\ref{tab:different_powers_of_i})

\begin{table}[h]
  \centering
  \caption{Different Powers of $i$}
  \label{tab:different_powers_of_i}

  \begin{tabular}{|l|c|c|r|}
    \hline
    $i = i$ & $i^{5} = i$ & $i^{9} = i$ & $i^{13} = i$ \\ \hline
    $i^{2} = -1$ & $i^{6} = -1$ & $i^{10} = -1$ & $i^{14} = -1$ \\ \hline
    $i^{3} = -i$ & $i^{7} = -i$ & $i^{11} = -i$ & $i^{15} = -i$ \\ \hline
    $i^{4} = 1$ & $i^{8} = 1$ & $i^{12} = 1$ & $i^{16} = 1$ \\ \hline
  \end{tabular}
\end{table}

\begin{example}
  \label{exm:different_powers_of_i}

  Simplify: $i^{27}$:

  First, you take the exponent ($27$) and divide it by $4$ and keep the
  remainder. The reason we divide it by $4$ is because that's how many version
  of $i$ there are.

  \longdivision[stage=1]{27}{4}

  We raise $i$ to the power of $3$, which is the remainder of $27 \div 4$.
  If we look at the different powers of $i$ table
  (\ref{tab:different_powers_of_i}), we will see that $i^{3} = -i$.
\end{example}

% subsubsection using_i_to_help_solve_equations (end)

\newpage
